\section{The free Maxwell stress and its collar Ward identities}
\label{sec:free-maxwell-ward-flux}

The reference stress is fixed by its action on every insertion, not only
by conservation. We prove that action using the same free curvature
covariance and Wick polynomial fields as in
Section~\ref{sec:ex-lattice-trace-free-rank}. All integrals involving
the stress below are taken on compact collars separated from every
insertion. No product of singular kernels at coincident points is used.

Let $\mathfrak g$ carry the fixed invariant real inner product. A smooth
$\mathfrak g$-valued two-form $f=(f_{ab})$ is a Maxwell wave on an
open set if it is antisymmetric and satisfies
\[
 \partial_af_{bc}+\partial_bf_{ca}+\partial_cf_{ab}=0,
 \qquad \sum_a\partial_af_{ab}=0.
\]
Indices run from zero to three. Set $\nu=(4\pi^2)^{-1}$ and
$G(z)=\nu|z|^{-2}$ for $z\ne0$. The covariance between curvature
components is
\begin{equation}\label{eq:maxwell-ward-covariance}
 C_{ab;il}(z)=\delta_{ib}\partial_a\partial_lG(z)
 -\delta_{lb}\partial_a\partial_iG(z)
 -\delta_{ia}\partial_b\partial_lG(z)
 +\delta_{la}\partial_b\partial_iG(z),
\end{equation}
times the identity on colour. A linear probe $p$ of a homogeneous
Maxwell jet of order $j$ acts on $D^jf(x)/j!$; write its value as
$p[f;x]$. Its covariance pulse $P_{p,x}(y)$ is the covariance of that
probe at $x$ with curvature at $y$. For a coordinate curvature probe
$(\gamma,i,l)$, with colour basis vector $e_\gamma$, this is $C_{ab;il}(y-x)e_\gamma$; higher pulses are obtained by
the corresponding normalized derivatives in $x$. Such probes span the
dual of the finite jet space, and these formulas extend linearly.
The raised probe $\partial_u p$ is defined by
$(\partial_up)[f;x]=\partial_u(p[f;\,\cdot\,])(x)$.
It includes the factor $j+1$ arising from normalized Taylor coefficients.
Differentiating the covariance gives
\begin{equation}\label{eq:maxwell-pulse-differentiation}
 \partial_{x,u}P_{p,x}(y)=P_{\partial_up,x}(y).
\end{equation}
Each pulse is a Maxwell wave away from its centre, as follows by
commuting derivatives in \eqref{eq:maxwell-ward-covariance} and using
$\Delta G=0$ there.

\subsection{The polarized current and its divergence}
For closed two-forms $f,g$, define the polarized stress
\[
 T_{ab}(f,g)=\sum_c\bigl(\langle f_{ac},g_{bc}\rangle
                   +\langle f_{bc},g_{ac}\rangle\bigr)
       -\frac{\delta_{ab}}2\sum_{c,d}\langle f_{cd},g_{cd}\rangle.
\]
It is symmetric in $a,b$ and in $f,g$, and its trace is zero in
four dimensions. Put $(\operatorname{div}f)_c=\sum_a\partial_af_{ac}$.
The Bianchi identity implies
\[
 \sum_{a,c}\langle f_{ac},\partial_ag_{bc}\rangle
       =\frac12\sum_{a,c}\langle f_{ac},\partial_bg_{ac}\rangle.
\]
Indeed substitute
$\partial_ag_{bc}=\partial_bg_{ac}-\partial_cg_{ab}$; exchanging
$a,c$ and using both antisymmetries makes the last sum equal to the
left side. The product rule, followed by this identity for $(f,g)$
and $(g,f)$, cancels the differentiated trace terms and gives
\begin{equation}\label{eq:maxwell-stress-divergence}
 \sum_a\partial_aT_{ab}(f,g)
 =\sum_c\bigl(\langle\operatorname{div}f_c,g_{bc}\rangle
                  +\langle f_{bc},\operatorname{div}g_c\rangle\bigr).
\end{equation}

We need $v(y)=e_b$ and $v(y)=A(y-x)$ for a real skew matrix $A$.
For either vector field define $j_a(v;f,g)=\sum_bv_bT_{ab}(f,g)$.
The term $\sum_{a,b}(\partial_av_b)T_{ab}$ vanishes by symmetry;
the same holds for dilation $v(y)=y-x$ by tracelessness.
Thus, if both forms are Maxwell waves, $\operatorname{div}j=0$.
For $\chi\in C_c^\infty(\mathbb R^4;\mathbb R)$ define outward
collar flux by
\begin{equation}\label{eq:maxwell-collar-flux}
          \mathcal F_\chi(j)=-\sum_a\int(\partial_a\chi)(y)j_a(y)\,dy.
\end{equation}
Whenever $j$ is smooth on a neighborhood of $\operatorname{supp}\chi$,
integration by parts identifies this with $\int\chi\operatorname{div}j$.
In particular it vanishes for two waves on that neighborhood.
If the only singularity is at $x$ and both cutoffs equal one near $x$,
their fluxes agree: apply the same identity to their difference, which
vanishes on a neighborhood of $x$.

\subsection{A unit-mass regulator fixes the sign}
For $t>0$ set
\[
 G_t(z)=\frac{\nu}{t^2+|z|^2},\qquad
 \rho_t(z)=\frac{8\nu t^2}{(t^2+|z|^2)^3}.
\]
Direct differentiation gives $-\Delta G_t=\rho_t$ and
$t^4\rho_t(tz)=\rho_1(z)$. These kernels are nonnegative, and polar
coordinates give
\[
 \int\rho_1(z)\,dz
 =16\pi^2\nu\int_0^\infty\frac{r^3}{(1+r^2)^3}\,dr=1.
\]
The radial integral is $1/4$, since an antiderivative is
$\tfrac14(1-(1+r^2)^{-1})^2$. For any continuous compactly supported
$h$, substitution $y=x+tz$ and dominated convergence therefore give
\begin{equation}\label{eq:maxwell-source-approximation}
              \int\rho_t(y-x)h(y)\,dy\longrightarrow h(x).
\end{equation}
The dominating function is $\rho_1(z)\|h\|_\infty$.

Fix a coordinate probe $(\gamma,i,l)$ and form a regulated pulse $R_t=dA_t$,
where
\[
 (A_t)_b(y)=\bigl(\delta_{ib}\partial_lG_t(y-x)
                    -\delta_{lb}\partial_iG_t(y-x)\bigr)e_\gamma.
\]
It is closed and satisfies
\[
 (\operatorname{div}R_t)_b
       =(\delta_{lb}\partial_i\rho_t(y-x)
                    -\delta_{ib}\partial_l\rho_t(y-x))e_\gamma.
\]
For the latter identity, $\sum_b\partial_b(A_t)_b=0$ by commuting
derivatives, so $\operatorname{div}dA_t=\Delta A_t$; now use
$\Delta G_t=-\rho_t$. Away from $x$, $R_t$ and all its derivatives
converge locally uniformly to $P_{(\gamma,i,l),x}$.

Suppose $f$ is a Maxwell wave on a neighborhood of
$\operatorname{supp}\chi$, and $\chi=1$ near $x$. Put
$a_b(y)=\sum_uv_u(y)f_{ub}(y)$ and
$H_b(y)=\chi(y)\langle e_\gamma,a_b(y)\rangle$.
Equation~\eqref{eq:maxwell-stress-divergence} gives
\[
 \mathcal F_\chi(j(v;R_t,f))
 =\int H_l\partial_i\rho_t(\,\cdot-x)
       -\int H_i\partial_l\rho_t(\,\cdot-x)
 =\int\rho_t(\,\cdot-x)\partial_lH_i
       -\int\rho_t(\,\cdot-x)\partial_iH_l.
\]
All integrations by parts are with smooth compactly supported $H_b$.
The right side tends to $\partial_lH_i(x)-\partial_iH_l(x)$ by
\eqref{eq:maxwell-source-approximation}. The left side tends to the
unregulated flux because $\operatorname{supp}d\chi$ is compact and
separated from $x$. Hence
\begin{equation}\label{eq:maxwell-coordinate-charge}
 \mathcal F_\chi(j(v;P_{(\gamma,i,l),x},f))
                       =-\langle e_\gamma,(d\iota_v f)_{il}(x)\rangle.
\end{equation}
This proves the sign directly from the Green normalization and the
outward-flux convention.

For an arbitrary jet probe $p$, differentiate this identity in the
enclosed centre while keeping both $v$ and $\chi$ fixed. The moving
centre stays in a neighborhood on which $\chi=1$. That neighborhood
can be chosen with positive distance from the compact collar. On its
product with the collar, the pulse derivatives in
\eqref{eq:maxwell-pulse-differentiation} are continuous and bounded.
They therefore provide an integrable constant times the collar indicator
for each fixed differentiation. Differentiation under the integral is
justified. Induction on the jet order and linearity in $p$ prove
\begin{equation}\label{eq:maxwell-single-charge}
        \mathcal F_\chi(j(v;P_{p,x},f))=-p[d\iota_v f;x].
\end{equation}
Normalized jets cause no change: the raising factors occur on both
sides and are nonzero. This argument applies whenever $d\iota_v f$
is a Maxwell wave locally, as it is for the affine motions used here.

For completeness, the Bianchi identity expands the variation as
\[
 (d\iota_v f)_{ab}
   =\sum_i\bigl(v_i\partial_if_{ab}
              +(\partial_av_i)f_{ib}+(\partial_bv_i)f_{ai}\bigr).
\]
For $v=e_b$ this is $\partial_bf$. For $v=A(y-x)$, it is the
infinitesimal inverse-rotation action on the jet at $x$. Differentiating
the two index factors and the argument in
$e^{tA}_{ia}e^{tA}_{lb}f_{il}(x+e^{tA}z)$ at $t=0$, where $z$
is displacement from the fixed rotation centre $x$, gives exactly
this formula. Orthogonal covariance preserves the Maxwell equations,
so this variation is a wave. Translation has the same property by
commutation of derivatives. Write $\mathscr D_vp$ for the resulting
probe: the raised derivative probe for translation and the angular
probe for rotation about $x$. Equation~\eqref{eq:maxwell-single-charge}
becomes $- (\mathscr D_vp)[f;x]$.

\subsection{Two contractions at the enclosed point}
When both pulses are centred at $x$, the preceding formula cannot be
used with one pulse as a smooth wave across $x$. Instead use collar
independence and homogeneity. For a probe of jet order $j$,
\[
 P_{p,x}(x+tz)=t^{-(4+j)}P_{p,x}(x+z)\qquad(t>0,z\ne0),
\]
because $G$ has degree $-2$ and the curvature covariance has two
derivatives, followed by $j$ centre derivatives. The drive has degree
$e=0$ for translation and $e=1$ for rotation about $x$.
Replace $\chi(y)$ by $\psi(y)=\chi(x+(y-x)/t)$. Both cutoffs
equal one near $x$, so their fluxes agree. In the latter integral set
$y=x+tz$. The volume contributes $t^4$, the cutoff derivative
$t^{-1}$, the drive $t^e$, and the two pulses
$t^{-8-j-l}$. Consequently their common flux $Z$ satisfies
\[
                   Z=t^{e-5-j-l}Z.
\]
Take $t=2$. Since $e\le1$ and $j,l\ge0$, the factor is strictly
less than one; hence $Z=0$. This proves vanishing without extending
the product of the two poles through $x$.

\subsection{All Wick polynomials and all spectators}
The Wick stress polynomial is
\begin{equation}\label{eq:maxwell-wick-stress}
 J^0_{ab}=\sum_c:\!\langle F_{ac},F_{bc}\rangle\!:
       -\frac{\delta_{ab}}4\sum_{c,d}:\!\langle F_{cd},F_{cd}\rangle\!:.
\end{equation}
Here the colons mean Wick ordering at that insertion. In Wick moments
there are no contractions between factors at the same insertion.
Fix distinct positions $x_0,\ldots,x_n$ and a smooth compact cutoff
$\chi$ equal to one near $x_0$, whose support excludes every other
$x_i$. Let $P_i$ be arbitrary Wick polynomials in finitely many jet
probes at $x_i$. Extend $\mathscr D_v$ from probes to polynomials
by the Leibniz rule, and denote this derivation by $\delta_v$. Here
$v=e_b$ for a coordinate direction $b$, or $v(y)=A(y-x_0)$ with $A$ real
skew; by linearity the claim then also holds for every affine Killing
field. We claim
\begin{equation}\label{eq:maxwell-polynomial-ward}
 -\sum_{a,b}\int(\partial_a\chi)(y)v_b(y)
   \left\langle J^0_{ab}(y)\prod_{i=0}^nP_i(x_i)\right\rangle_0dy
       =-\left\langle(\delta_vP_0)(x_0)
                            \prod_{i=1}^nP_i(x_i)\right\rangle_0.
\end{equation}

First suppose each $P_i$ is a monomial. Label its individual factor
occurrences, including repeated factors, by a finite set $I$. Let
$x_{u(h)}$ be the position of occurrence $h$, with probe $p_h$,
and put $f_h=P_{p_h,x_{u(h)}}$. For a subset of remaining occurrences
$A$, write $W(A)$ for its Wick pairing sum, with pair weights zero
at the same insertion. Extract the two stress factors from the Wick
sum. The exact integrand contracted with $v_b$ is
\begin{equation}\label{eq:maxwell-double-extraction}
 \frac12\sum_{h\ne l}j_a(v;f_h,f_l)(y)\,
                                      W(I\setminus\{h,l\}).
\end{equation}
To verify the factor $1/2$, initially assign the first stress factor
to $h$ and the second to $l$. The quadratic expression in
\eqref{eq:maxwell-wick-stress} gives one ordered contraction.
Exchanging $h,l$ gives its transpose; their sum is exactly the
polarized tensor $T_{ab}(f_h,f_l)$. Averaging these two equal
ordered sums gives \eqref{eq:maxwell-double-extraction}.
The stress factors cannot contract to each other because of Wick
ordering. After choosing their partners, the remaining pairing sum
is exactly $W(I\setminus\{h,l\})$.

Every term is smooth on the compact collar, so the finite sum can be
integrated term by term. There are three cases. If neither partner
is at $x_0$, both pulses are smooth waves on the support of $\chi$,
and the flux vanishes. If both are at $x_0$, the preceding homogeneous
two-pulse argument gives zero. If exactly one, say $h$, is at $x_0$,
\eqref{eq:maxwell-single-charge} gives minus the covariance of
$\mathscr D_vp_h$ at $x_0$ with $p_l$ at its spectator position.
Both orders $(h,l)$ occur, and symmetry cancels the factor $1/2$.
For a fixed $h$ at $x_0$, sum over its possible partner $l$ and
over the remaining Wick pairings. This is precisely the Wick moment
with occurrence $h$ replaced by $\mathscr D_vp_h$. Summing over all
such occurrences is the Leibniz rule for $\delta_vP_0$, proving
\eqref{eq:maxwell-polynomial-ward} for monomials. Expand general
polynomials in their finite monomial bases. Multilinearity and the
linearity of the compact integrals prove the identity for arbitrary
complex coefficients and every finite spectator word.

\subsection{The fixed reference normalization}
For this final step assume the gauge Lie algebra is nonzero, nonabelian
and simple, as in the construction theorem. The preceding collar and Wick
calculations used only its finite-dimensional real inner product; the
rank and nontriviality assertions below use this additional hypothesis.
The polynomial \eqref{eq:maxwell-wick-stress} is real, symmetric,
gauge-invariant and of grade four. Its expectation vanishes by Wick
ordering. Its divergence is zero already as a Maxwell jet polynomial:
apply \eqref{eq:maxwell-stress-divergence} with $f=g$ and divide
by two. Polynomial differentiation is the same Leibniz rule as the
field derivative. Its reference field is nonzero by the explicit
stress probes in Section~\ref{sec:ex-lattice-trace-free-rank}.
That section also proves surjectivity of divergence on centered
covariant arrays and that its kernel is the line through this tensor.
These facts supply the conserved-line certificate used in
Section~\ref{sec:actual-stress-normalization}.

We use three properties of the reference field table. They are part of
its construction and are stated inputs of this section; they are not proved
here.
\begin{enumerate}
\item[(R1)] Faithfulness: a field label whose separated correlations with
  every spectator word vanish is zero.
\item[(R2)] The reference field labels are the images of the invariant
  Wick jet polynomials, and this identification preserves Wick moments,
  sends the raised probe $\partial_up$ to the derivative plan in direction
  $u$, and sends the rotation variation $\mathscr D_vp$, $v(y)=A(y-x)$, to
  the infinitesimal rotation action.
\item[(R3)] The canonical reference stress $J^0$ is the image of the Wick
  polynomial \eqref{eq:maxwell-wick-stress}.
\end{enumerate}
By (R2) and (R3), \eqref{eq:maxwell-polynomial-ward} descends to all field
labels. For the reference tensor $\tau J^0$, $\tau\in\mathbb R
\setminus\{0\}$, both its translation and rotation Ward maps have
the common coefficient $-\tau$. The sign comes from the collar
calculation, not a choice made after comparing with the physical theory.

There is also a useful uniqueness check on the reference packaging.
Suppose the same stress field has translation Ward identities with
coefficients $c_0,c_1$ multiplying the same derivative plan. Subtracting
the two identities shows that every separated correlation with
$(c_0-c_1)\partial_bv$ vanishes. Faithfulness (R1)
gives $(c_0-c_1)\partial_bv=0$ for every label $v$. The derivative
plan is nonzero already in grade four, as detected by the nonzero
divergence probes in Section~\ref{sec:ex-lattice-trace-free-rank}; that
detection uses the compatibility of the derivative plan with the field
derivatives in (R2).
Hence $c_0=c_1$. In particular, retaining the canonical tensor $J^0$
and the same derivative plan forces the coefficient to be $-1$.

In the native construction the reference is fixed before the weak
comparison curve is chosen. By (R3) its canonical tensor has polynomial
\eqref{eq:maxwell-wick-stress}: the quadratic coordinate labels map
to the corresponding products of curvature probes, their sum gives
the first term, and subtracting one quarter of the diagonal sum
gives the second. Selecting this tensor does not change the
hierarchy, derivative plan, rotation or vacuum of the reference table, so
(R1) and (R2) apply to the same data.
The uniqueness argument therefore fixes its common Ward coefficient
to $-1$ on that same data. Combined with the actual conserved-line
certificate, this justifies the reference input to the physical
normalization step without substituting another comparison model.
